\title{Towards a proteome-scale map of the human protein-protein interaction network}
\begin{document}
\maketitle

\begin{abstract}
© 2005 Nature Publishing Group Towards a proteome-scale map of the human protein–protein interaction network Jean-Franc¸ois Rual1*, Kavitha Venkatesan1*, Tong Hao1, Tomoko Hirozane-Kishikawa1, Ame´lie Dricot1, Ning Li1, Gabriel F. Berriz2, Francis D. Gibbons2, Matija Dreze1,3, Nono Ayivi-Guedehoussou1, Niels Klitgord1, Christophe Simon1, Mike Boxem1, Stuart Milstein1, Jennifer Rosenberg1, Debra S. Goldberg2, Lan V. Zhang2, Sharyl L. Wong2, Giovanni Franklin2, Siming Li1†, Joanna S. Albala1†, Janghoo Lim4, Carlene Fraughton1, Estelle Llamosas1, Sebiha Cevik1, Camille Bex1, Philippe Lamesch1,3, Robert S. Sikorski5, Jean Vandenhaute3, Huda Y. Zoghbi4, Alex Smolyar1, Stephanie Bosak6, Reynaldo Sequerra6, Lynn Doucette-Stamm6, Michael E. Cusick1, David E. Hill1, Frederick P. Roth2 \& Marc Vidal1 Systematic mapping of protein–protein interactions, or ‘inter- actome’ mapping, was initiated in model organisms, starting with defined biological processes1,2 and then expanding to the scale of the proteome3–7. Although far from complete, such maps have revealed global topological and dynamic features of interactome networks that relate to known biological properties8,9, suggesting that a human interactome map will provide insight into develop- ment and disease mechanisms at a systems level. Here we describe an initial version of a proteome-scale map of human binary protein–protein interactions. Using a stringent, high-throughput yeast two-hybrid system, we tested pairwise interactions amon
\end{abstract}

\section{Recovered Text}
© 2005 Nature Publishing Group
Towards a proteome-scale map of the human
protein–protein interaction network
Jean-Franc¸ois Rual1*, Kavitha Venkatesan1*, Tong Hao1, Tomoko Hirozane-Kishikawa1, Ame´lie Dricot1, Ning Li1,
Gabriel F. Berriz2, Francis D. Gibbons2, Matija Dreze1,3, Nono Ayivi-Guedehoussou1, Niels Klitgord1,
Christophe Simon1, Mike Boxem1, Stuart Milstein1, Jennifer Rosenberg1, Debra S. Goldberg2, Lan V. Zhang2,
Sharyl L. Wong2, Giovanni Franklin2, Siming Li1†, Joanna S. Albala1†, Janghoo Lim4, Carlene Fraughton1,
Estelle Llamosas1, Sebiha Cevik1, Camille Bex1, Philippe Lamesch1,3, Robert S. Sikorski5, Jean Vandenhaute3,
Huda Y. Zoghbi4, Alex Smolyar1, Stephanie Bosak6, Reynaldo Sequerra6, Lynn Doucette-Stamm6,
Michael E. Cusick1, David E. Hill1, Frederick P. Roth2 \& Marc Vidal1
Systematic mapping of protein–protein interactions, or ‘inter-
actome’ mapping, was initiated in model organisms, starting with
defined biological processes1,2 and then expanding to the scale of
the proteome3–7. Although far from complete, such maps have
revealed global topological and dynamic features of interactome
networks that relate to known biological properties8,9, suggesting
that a human interactome map will provide insight into develop-
ment and disease mechanisms at a systems level. Here we describe
an initial version of a proteome-scale map of human binary
protein–protein interactions. Using a stringent, high-throughput
yeast two-hybrid system, we tested pairwise interactions among
the products of ,8,100 currently available Gateway-cloned open
reading frames and detected ,2,800 interactions. This data set,
called CCSB-HI1, has a verification rate of ,78\% as revealed by an
independent co-affinity purification assay, and correlates signifi-
cantly with other biological attributes. The CCSB-HI1 data set
increases by ,70\% the set of available binary interactions within
the tested space and reveals more than 300 new connections to
over 100 disease-associated proteins. This work represents an
important step towards a systematic and comprehensive human
interactome project.
Our working definition of a human interactome map is the
complete collection of binary protein–protein interactions detectable
in one or more exogenous assay. This definition excludes dynamic
and functional properties of these interactions (Supplementary
Data I). Thus, we treat interactome maps as ‘scaffold’ information,
from which increasingly detailed and reliable biological models can
be generated by integrating other functional genomic and proteomic
data sets10 (Supplementary Data II).
The currently available information on the human interactome
network originates from either literature-curated (LC) inter-
actions11–15, or from ‘interologs’ (that is, potential interactions
predicted from interactome data available for model organisms
given evolutionary conservation of two known partners)2,16. This
information needs to be complemented by systematic experimental
mapping approaches that are: (1) not biased towards any particular
biological interest (that is, without ‘inspection bias’), as is the case for
LC data sets; (2) more complete; and (3) supported by experiments
rather than predictions. We are mapping the human interactome
network systematically in successive versions, with each version
defined by the availability of recombinationally cloned open reading
frames (ORFs) in the human ‘ORFeome’17.
In this initial version, we use human ORFeome v1.1 (ref. 17), a
resource containing ,8,100 Gateway-cloned ORFs (generated using
the Mammalian Gene Collection, as previously described17) that
correspond to ,7,200 distinct protein-coding genes (Supplementary
Table S1). Thus, our initial ‘search space’ (Space-I) encompasses
protein pairs encoded by a 7,200 £ 7,200 matrix of genes. Future
interactome versions can be generated by successively increasing the
search space as additional versions of the human ORFeome become
available (Supplementary Data III). Accepting a total of ,22,000
protein-coding genes in the human genome18 and excluding poly-
morphic and splice variants, Space-I corresponds to ,10\% of the
total search space for a comprehensive human interactome map
(Fig. 1a). Currently, 4,067 binary LC interactions are available in
Space-I (LCI interactions; Supplementary Table S2).
Our high-throughput yeast two-hybrid system is highly specific,
benefiting from the following features, which were not uniformly
present in earlier large-scale studies: relatively low levels of expression
of both Gal4 DNA-binding domain (DB) and Gal4 activation
domain (AD) hybrid proteins (DB-X and AD-Y, or DB-ORF and
AD-ORF); three different yeast two-hybrid inducible reporter genes;
and a plasmid-shuffling counter selection to eliminate systematically
de novo auto-activators19 (Supplementary Data IV). We tested each of
the ,8,100 individual DB-X proteins against 45 mini-libraries, each
containing a pool of 188 AD-Y fusion proteins (AD-188Ys), by
yeast-mating in a 96-well format (Fig. 1a). Such small pools
offer high sensitivity, because positive clones are less likely to be
masked by other AD-Y clones within the same pool. Indeed, our
overall reproducibility rate was ,55\%, close to that observed in
proteome-scale affinity purification followed by mass spectrometry
experiments20 (Supplementary Data V).
In our Space-I yeast two-hybrid matrix, we identified ,65,000
primary positive colonies, of which 12,251 scored positive after
LETTERS
1Center for Cancer Systems Biology and Department of Cancer Biology, Dana-Farber Cancer Institute and Department of Genetics, Harvard Medical School, 44 Binney Street,
Boston, Massachusetts 02115, USA. 2Department of Biological Chemistry and Molecular Pharmacology, Harvard Medical School, 250 Longwood Ave, Boston, Massachusetts
02115, USA. 3Unite´ de Recherche en Biologie Mole´culaire, Faculte´s Notre-Dame de la Paix, 61 Rue de Bruxelles, 5000 Namur, Belgium. 4Howard Hughes Medical Institute, and
Departments of Pediatrics, Neurology, Neuroscience, and Molecular and Human Genetics, Baylor College of Medicine, One Baylor Plaza, Houston, Texas 77030, USA. 5Arcbay,
Inc., 6 Whittier Place, Suite 7J, Boston, Massachusetts 01915, USA. 6Agencourt Bioscience Corporation, 500 Cummings Center, Suite 2450, Beverly, Massachusetts 01915, USA.
†Present addresses: ArQule, Inc., 19 Presidential Way, Woburn, Massachusetts 01081, USA (S.L.); Departments of Cancer Biology, and Otolaryngology, Head and Neck Surgery,
University of California Davis, 2521 Stockton Blvd, Suite 7200, Sacramento, California 95817, USA (J.S.A.).
*These authors contributed equally to this work.
Vol 437|20 October 2005|doi:10.1038/nature04209
1173
© 2005 Nature Publishing Group
stringent phenotype testing; that is, we only retained clones that were
positive for at least two yeast two-hybrid reporter assays and we
controlled for auto-activation (Fig. 1a). Both DB-X and AD-Y
fragments from these two-reporter-positive colonies were amplified
by polymerase chain reaction (PCR) and sequenced to generate
10,597 pairs of interaction sequence tags (ISTs) (Supplementary
Fig. S1a). We then collapsed yeast two-hybrid ISTs corresponding
to the same pair of genes and removed lower confidence interactions
(Supplementary Data VI). This resulted in a data set containing 2,754
yeast two-hybrid interactions—the Center for Cancer Systems
Biology Human Interactome version 1 (CCSB-HI1) (Supplementary
Fig. S1a and Supplementary Table S2. All data is publicly available
(Supplementary Data VII). In all, CCSB-HI1 provides interaction
information for 1,549 proteins, ,21\% of the proteins tested in
Space-I.
To measure the specificity of CCSB-HI1, we considered technical
and biological false positives separately (Supplementary Data VIII).
Technical false positives arise from experimental errors that can and
should be avoided. To estimate our technical false positive rate,
representative samples of interactions were verified by in vivo co-
affinity purification glutathione S-transferase (GST) pull-down assay
in human 293T cells (Supplementary Data IX). The verification rates
for co-affinity purification were ,78\% for yeast two-hybrid-only
interactions, ,62\% for LCI-only interactions and ,81\% for inter-
actions present in both LCI and yeast two-hybrid data sets (Fig. 1b; see
also Supplementary Fig. S1b, c and Supplementary Tables S2 and S3).
Given these results and the fact that co-affinity purification GST pull-
down assays are not perfectly sensitive, we argue that CCSB-HI1 is
largely free of technical false positives, and comparable in reliability
to LCI interactions, supporting the validity of our improvements of
Figure 1 | Towards the generation of a proteome-scale human yeast two-
hybrid map. a, Schema of the high-throughput yeast two-hybrid pipeline.
Individual steps (middle column) and representative examples (flanking left
and right columns) are indicated. The top panel of the left column
represents the matrix of all protein pairs. All available ORFs from human
ORFeome v1.1 were transferred into both DB and AD vectors by
recombinational cloning (middle panel of left column). The top panel of the
right column shows the mating process, with each bait mated to individual
pools of 188 AD-ORFs. Initial phenotypic testing evaluated growth of
diploid cells on selective medium in response to enhanced levels of the
GAL1::HIS3 selective marker (bottom panel of left column). All positive
diploids from phenotyping no. 1 (red circles) were subsequently tested for
activation of both GAL1::HIS3 and GAL1::lacZ reporter genes. Auto-
activators were identified by growth on medium containing cycloheximide
(bottom panels of left and right columns). Positive colonies from
phenotyping no. 2 (outlined in red) were isolated and used to PCR-amplify
both DB-ORF and AD-ORF fragments for sequencing. b, Verification of
yeast two-hybrid interactions by co-affinity purification assays. Fifteen
representative examples of co-affinity purification-positive assays are
shown. The middle and bottom panels show expression controls of
Myc–prey and GST–bait fusion proteins, respectively. Each lane pair in the
top panels shows presence or absence of Myc–prey fusions after affinity
purification, demonstrating binding to GST–bait fusion proteins (þ) or to
GST alone (2). The Table summarizes the data obtained for four different
classes of protein pairs. ‘Y2H and LCI’ describes interactions reported in
both the yeast two-hybrid and LCI data sets. ‘Y2H/LCI-negative’ describes
pairs of proteins that were not reported to interact either in the yeast two-
hybrid or in the LCI data sets. Rows indicate the total number of interactions
tested and considered for scoring (Total), the number of interactions not
verified by co-affinity purification (co-AP2), the number of interactions
verified by co-affinity purification (co-APþ), the proportion of co-affinity
purification-positive interactions (success rate), and the adjusted success
rate (which accounts for the observation that one-third of all co-affinity
purification experiments yield an apparently positive result without regard
to whether or not the protein pair truly interacts; see Supplementary Data
IX). Identities, lane positions and scoring of all protein pairs tested by
co-affinity purification are provided in Supplementary Tables S2 and S3.
LETTERS
NATURE|Vol 437|20 October 2005
1174
© 2005 Nature Publishing Group
the yeast two-hybrid methodology. Estimating biological false posi-
tive interactions, which are genuinely observed in one or more assay
but do not occur in vivo, is more difficult. We partially addressed this
by examining the correlation of CCSB-HI1 data with other biological
information (see below).
To measure the sensitivity of CCSB-HI1, we selected two high-
confidence subsets from among all 4,067 LCI direct binary inter-
actions. LCI-core contains 624 interactions supported by at least two
PubMed entries. LCI-hypercore contains 275 interactions supported
by at least two PubMed entries and present in at least two curated
databases (Supplementary Table S2). Overall, the fractions of LCI, LCI-
core and LCI-hypercore interactions found in CCSB-HI1 are 2.3\%,
4.6\% and 8.4\%, respectively (Fig. 2a). These overlaps are larger than
expected by chance (P , 6 £ 10256) and are similar to those found
for interactome maps in Caenorhabditis elegans and Drosophila
melanogaster7,21. That the fraction of CCSB-HI1 interactions increases
markedly with increasingly confident subsets of LCI suggests that
literature-derived interactions are variable in quality and should not
necessarily be interpreted as a ‘gold standard’. Because Space-I rep-
resents ,10\% of the human network (without accounting for
alternative splice variants), and because we detected ,10\% of LCI-
hypercore interactions, we conclude that the CCSB-HI1 data set
contains ,1\% of the human interactome (Supplementary Data X).
We represented the union of all CCSB-HI1 and LCI interactions in
a network graph in which nodes are proteins and edges are inter-
actions. The main component of this network contains 2,784 nodes
and 6,438 edges (Fig. 2b), and shows interactions largely segregated
into two neighbourhoods: one enriched for CCSB-HI1 interactions
(red edges) and the other enriched for LCI interactions (blue edges).
To explore this hypothesis, we calculated, for each node, the fraction
of yeast two-hybrid edges within paths of length 1, 2 and 3 (that is,
within ‘1-hop’, ‘2-hop’ and ‘3-hop’ neighbourhoods). The distri-
bution of this fraction (Fig. 2c; see also Supplementary Fig. S2)
confirms the evidence-type segregation apparent in Fig. 2b. One
explanation for this phenomenon is that different biases exist in the
CCSB-HI1 and LCI data sets. For example, certain protein classes
(such as those involved in cancer) are studied more extensively
than others, resulting in an inherent inspection bias in LC data
(Supplementary Table S4). Furthermore, the methodologies used to
detect interactions (including yeast two-hybrid) each have different
biases for example, under-representation of membrane proteins
(Supplementary Data X).
The novelty of CCSB-HI1 interactions was evaluated by system-
atically searching the PubMed and Google Scholar literature data-
bases for co-occurrence of the corresponding gene symbols. More
than 85\% of the CCSB-HI1 pairs (as compared with only 25\% of
Figure 2 | Overlap of CCSB-HI1 with existing literature-curated (LC)
data. a, Overlap between CCSB-HI1 and LC interactions in Space-I (LCI).
The top, middle and bottom panels represent the overlap between
CCSB-HI1 and LCI, LCI-core and LCI-hypercore, respectively. b, Network
graph of the union of all CCSB-HI1 and LCI interactions. Proteins are shown
as yellow nodes and CCSB-HI1 and LCI interactions are shown as red and
blue edges, respectively. Blue edges with increasing thickness indicate
LCI-non-core, LCI-core and LCI-hypercore, respectively. The apparent
banding pattern of the yellow nodes is an artefact of the graph layout
algorithm (Supplementary Data). Importantly, the layout algorithm was not
informed by type of supporting evidence and therefore does not explain the
evident separation of blue and red edges. c, Bias in 2-hop network
neighbourhood for either CCSB-HI1 or LCI interactions. The frequency of
nodes with a given proportion of CCSB-HI1 interactions in their 2-hop
neighbourhood is depicted for the interactome network graph in b (solid
curve) and for a network in which the types of supporting evidence (CCSB-
HI1 or LCI) are randomly permuted among edges (dashed curve). The solid
curve indicates that most of the proteins in the network of b have either only
CCSB-HI1 or only LCI interactions in their 2-hop neighbourhood. In
contrast, neighbourhoods are well mixed when evidence labels are randomly
permuted among edges.
NATURE|Vol 437|20 October 2005
LETTERS
1175
© 2005 Nature Publishing Group
Figure 3 | Interaction network of disease-associated CCSB-HI1
proteins. The network has 121 OMIM disease-associated proteins (green
nodes) and 424 CCSB-HI1 interactions involving them (red edges), along
with known LC interactions (solid blue edges represent binary LCI
interactions and dashed blue edges represent non-binary interactions).
Proteins without an OMIM disease association are depicted as yellow nodes,
and blue edges with increasing thickness indicate LCI-non-core, LCI-core
and LCI-hypercore interactions, respectively. We note that 94 out of the 424
CCSB-HI1 interactions involve the Ewing sarcoma related protein (EWSR1;
also known as EWS).
Table 1 | Overlap of protein interactions with other gene- or protein-pair characteristics
Protein pairs
Share mouse phenotype
Share upstream motif
Have correlated expression
F(C)
P-value
F(C)
P-value
F(C)
P-value
All possible within Space-I
0.128
NA
0.086
NA
0.063
NA
CCSB-HI1
0.257
2.53 £ 1023
0.115
1.14 £ 1024
0.130
2.14 £ 1027
LCI
0.336
4.91 £ 10243
0.146
9.05 £ 10220
0.204
5.45 £ 10256
LCI-core
0.471
7.53 £ 10220
0.137
3.77 £ 1023
0.243
3.57 £ 10212
LCI-non-core
0.306
2.54 £ 10227
0.147
3.65 £ 10218
0.198
7.78 £ 10246
Protein pairs
Share GO component
Share GO function
Share GO process
F(C)
P-value
F(C)
P-value
F(C)
P-value
All possible within Space-I
0.059
NA
0.021
NA
0.036
NA
CCSB-HI1
0.488
1.49 £ 10228
0.250
5.49 £ 10220
0.233
2.68 £ 10228
LCI
0.656
6.47 £ 102139
0.228
5.72 £ 102120
0.410
1.74 £ 102405
LCI-core
0.870
5.90 £ 10243
0.270
1.70 £ 10232
0.616
3.40 £ 102137
LCI-non-core
0.610
3.84 £ 102100
0.218
1.33 £ 10289
0.368
4.03 £ 102280
F(C) represents the fraction of gene- or protein-pairs (defined for each row) the given characteristic C. Assessed characteristics include shared mouse phenotype, shared upstream motif,
correlated expression30 and shared Gene Ontology annotation. ‘All possible within Space-I’ represents all possible gene- or protein-pairs in Space-I for which information regarding C is
available. For each analysis of a shared characteristic, only gene- or protein-pairs for which both members had some annotation for that characteristic were considered. For analysis of
correlated expression, only gene pairs with expression measurements for both genes were considered.
LETTERS
NATURE|Vol 437|20 October 2005
1176
© 2005 Nature Publishing Group
LCI pairs) showed no linkage of the corresponding gene symbols
(Supplementary Fig. S3). These results indicate that most of our yeast
two-hybrid interactions are novel.
To determine whether messenger RNAs corresponding to inter-
acting protein pairs are likely to be co-expressed, we used Pearson
correlation coefficients of the corresponding gene pairs in the
CCSB-HI1 and LCI data sets from four expression studies in diverse
human and mouse tissues (Supplementary Data XI). LCI pairs were
enriched for correlated expression in all four cases (P , 3 £ 10217)
and CCSB-HI1 pairs were enriched in three of the four cases
(P , 3 £ 1025) (Table 1; see also Supplementary Fig. S4a and
Supplementary Tables S2 and S5). In addition, CCSB-HI1 inter-
actions are more enriched than would be expected by chance for (1)
presence of a common upstream DNA sequence that is conserved
across human, mouse, rat and dog genomes22 (P ¼ 1 £ 1024), (2)
orthologous genes in mouse having a specific phenotype in common23
(P ¼ 3 £ 1023), and (3) annotation with the same Gene Ontology
(GO) terms24 (P , 6 £ 10220 for all three GO branches) (Table 1; see
also Supplementary Tables S2 and S5 and Supplementary Data XI).
The higher likelihood of LCI interactions to share other biological
attributes is not surprising given inspection bias and potential
circularity where functional annotation has been derived from an
LCI interaction. That the CCSB-HI1 interaction pairs (which do not
have such biases) yield statistically significant correlation supports
their biological relevance. In all, 357 CCSB-HI1 interactions are
supported by at least one additional characteristic and thus represent
particularly appealing hypotheses of functional relatedness (Sup-
plementary Fig. S4b). Lack of additional biological evidence is not an
argument against any interaction (Supplementary Data XII). Impor-
tantly, complementary information from the interactome and other
functional genomic data can be integrated to formulate biological
models25.
The CCSB-HI1 network has topological properties that are similar
to other sampled interactome networks, such as an approximately
power-law degree distribution, hierarchical organization and a
tendency for highly connected (hub) proteins to interact with
less highly connected proteins (Supplementary Data XIII and Sup-
plementary Fig. S5a–d). Surprisingly, although the CCSB-HI1 net-
work has a small characteristic path length, it does not exhibit high
clustering, seemingly contradicting findings from the sparsely
sampled networks of other organisms that protein interaction net-
works are ‘small world’6,7,26 (that is, have a short characteristic path
length and a high clustering coefficient). Possible explanations for
this apparent discrepancy are discussed in Supplementary Data XIII.
To gain an insight into the evolution of the interactome, we
classified proteins in the CCSB-HI1 network as ‘eukaryotic’,
‘metazoan’, ‘mammalian’ or ‘human’, and asked whether proteins
specific to different evolutionary classes tend to interact with one
another. The CCSB-HI1 network appears to be enriched for inter-
actions between proteins of the same evolutionary class but not for
interactions between proteins from two different evolutionary classes
(Supplementary Table S6). This suggests that the human interactome
has evolved through the preferential addition of interactions
between lineage-specific proteins. Further investigation may provide
hypotheses for mechanisms underlying interactome evolution.
To detect densely connected subgraphs potentially representing
biological modules, we applied the MCODE graph clustering algo-
rithm27 to the CCSB-HI1 and to the combined CCSB-HI1–LCI and
CCSB-HI1–LC networks (Supplementary Fig. S6, Supplementary
Table S7 and Supplementary Data XIV). We identified functionally
enriched MCODE complexes using FuncAssociate28. Out of 172
complexes (Supplementary Table S7) containing at least one
CCSB-HI1 interaction, we identified 102 in which at least one GO
term was significantly over-represented (P \# 0.05), ten times the
number expected by chance alone. The enriched functional terms we
identified may also apply to unannotated proteins present in the
complex (‘guilt by association’ predictions).
CCSB-HI1 represents a repository of novel biological hypotheses
for genes implicated in human diseases. We compared all CCSB-HI1
proteins to the list of genes associated with human diseases in
the Online Mendelian Inheritance in Man (OMIM) database and
identified 424 interacting pairs for which at least one partner had
been previously associated with a human disease (Supplementary
Table S8). In a query of PubMed and Google Scholar, searching for
gene symbols, 352 of the 424 interaction pairs appeared to be new
based on the absence of any hit in either database. Along with 79 LC
interactions (including LCI and non-binary LC interactions) among
proteins in this space, the resulting network contains 484 interactions
among 417 proteins (Fig. 3).
In one example, we found an interaction between RTN4, a neurite
outgrowth inhibitor, and SPG21, the spastic paraplegia 21 protein.
Mutations in SPG21 cause an autosomal recessive motor disorder
called Mast syndrome. SPG21 protein localizes to intracellular
endosomal/trans-Golgi transportation vesicles and is thought to
function in protein transport and sorting. Although the function
of its interacting partner, RNT4, remains elusive, RNT4 belongs to a
family of proteins that localize to the endoplasmic reticulum and are
markers of neuroendocrine differentiation29. In addition, RNT4
shares two regulatory motifs with SPG21 that are conserved across
mammalian genomes22 and may have a role in Mast syndrome. This
and other examples (Supplementary Data XV) suggest that CCSB-
HI1 can be used to connect biological processes in order to under-
stand further network and disease relationships.
Although the CCSB-HI1 data set is far from comprehensive, and
incomplete sampling limits conclusions regarding some network
properties, our data provide useful hypotheses and can guide further
studies of the expanded network. Currently, CCSB-HI1 is a static
graph. Eventually the dynamics of the human interactome network
will need to be considered to address where and when interactions
take place and how they are regulated. The functional consequences
of these physical interactions will also have to be studied to under-
stand the logic of complex biological networks. Just as the first drafts
of the human genome changed strategies for disease gene identifi-
cation, the emerging human interactome will greatly further the
understanding of human health and disease.
METHODS
The CCSB-HI1 data set was generated using a high-throughput version of the
yeast two-hybrid system. First, all 8,100 cloned ORFs of the human ORFeome
v1.1 were transferred from Entry clones to both AD and DB vectors by
Gateway recombinational cloning. Resulting constructs were transformed in
haploid yeast cells. After mating, diploids were tested on selective media for
their ability to grow in an AD-Y-dependent manner. The identity of the
interactors was determined after PCR amplification and sequencing of the AD
and DB inserts from positive colonies. Resulting interactions were re-tested by
the yeast two-hybrid system, individually, assessed for quality by co-affinity
purification assays and analysed for correlation with other biological infor-
mation. For a detailed description of the various methods, see Supplementary
Data.
Received 21 July; accepted 8 September 2005.
Published online 28 September 2005.
1.
Fromont-Racine, M., Rain, J. C. \& Legrain, P. Toward a functional analysis of
the yeast genome through exhaustive two-hybrid screens. Nature Genet. 16,
277–-282 (1997).
2.
Walhout, A. J. et al. Protein interaction mapping in C. elegans using proteins
involved in vulval development. Science 287, 116–-122 (2000).
3.
Uetz, P. et al. A comprehensive analysis of protein-protein interactions in
Saccharomyces cerevisiae. Nature 403, 623–-627 (2000).
4.
Ito, T. et al. A comprehensive two-hybrid analysis to explore the yeast protein
interactome. Proc. Natl Acad. Sci. USA 98, 4569–-4574 (2001).
5.
Reboul, J. et al. C. elegans ORFeome version 1.1: experimental verification of the
genome annotation and resource for proteome-scale protein expression.
Nature Genet. 34, 35–-41 (2003).
6.
Giot, L. et al. A protein interaction map of Drosophila melanogaster. Science 302,
1727–-1736 (2003).
7.
Li, S. et al. A map of the interactome network of the metazoan C. elegans.
Science 303, 540–-543 (2004).
NATURE|Vol 437|20 October 2005
LETTERS
1177
© 2005 Nature Publishing Group
8.
Jeong, H., Mason, S. P., Barabasi, A. L. \& Oltvai, Z. N. Lethality and centrality in
protein networks. Nature 411, 41–-42 (2001).
9.
Han, J. D. et al. Evidence for dynamically organized modularity in the yeast
protein–-protein interaction network. Nature 430, 88–-93 (2004).
10. Vidal, M. A biological atlas of functional maps. Cell 104, 333–-339 (2001).
11.
Xenarios, I. et al. DIP, the Database of Interacting Proteins: a research tool for
studying cellular networks of protein interactions. Nucleic Acids Res. 30,
303–-305 (2002).
12.
Zanzoni, A. et al. MINT: a Molecular INTeraction database. FEBS Lett. 513,
135–-140 (2002).
13.
Bader, G. D., Betel, D. \& Hogue, C. W. BIND: the Biomolecular Interaction
Network Database. Nucleic Acids Res. 31, 248–-250 (2003).
14.
Peri, S. et al. Development of human protein reference database as an initial
platform for approaching systems biology in humans. Genome Res. 13,
2363–-2371 (2003).
15.
Pagel, P. et al. The MIPS mammalian protein-protein interaction database.
Bioinformatics 21, 832–-834 (2005).
16.
Lehner, B. \& Fraser, A. A first-draft human protein-interaction map. Genome
Biol. 5, R63 (2004).
17.
Rual, J. F. et al. Human ORFeome version 1.1: a platform for reverse proteomics.
Genome Res. 14, 2128–-2135 (2004).
18.
International Human Genome Sequencing Consortium. Finishing the
euchromatic sequence of the human genome. Nature 431, 931–-945 (2004).
19.
Walhout, A. J. \& Vidal, M. A genetic strategy to eliminate self-activator baits
prior to high-throughput yeast two-hybrid screens. Genome Res. 9, 1128–-1134
(1999).
20. Gavin, A. C. et al. Functional organization of the yeast proteome by systematic
analysis of protein complexes. Nature 415, 141–-147 (2002).
21.
Formstecher, E. et al. Protein interaction mapping: A Drosophila case study.
Genome Res. 15, 376–-384 (2005).
22. Xie, X. et al. Systematic discovery of regulatory motifs in human promoters
and 3
0 UTRs by comparison of several mammals. Nature 434, 338–-345
(2005).
23. Eppig, J. T. et al. The Mouse Genome Database (MGD): from genes to mice—a
community resource for mouse biology. Nucleic Acids Res. 33, D471–-D475
(2005).
24. Camon, E. et al. The Gene Ontology Annotation (GOA) Database: sharing
knowledge in Uniprot with Gene Ontology. Nucleic Acids Res. 32, D262–-D266
(2004).
25. Gunsalus, K. C. et al. Predictive models of molecular machines involved in
Caenorhabditis elegans early embryogenesis. Nature 436, 861–-865 (2005).
26. Wagner, A. The yeast protein interaction network evolves rapidly and contains
few redundant duplicate genes. Mol. Biol. Evol. 18, 1283–-1292 (2001).
27. Bader, G. D. \& Hogue, C. W. An automated method for finding molecular
complexes in large protein interaction networks. BMC Bioinformatics 4, 2
(2003).
28. Berriz, G. F., King, O. D., Bryant, B., Sander, C. \& Roth, F. P. Characterizing gene
sets with FuncAssociate. Bioinformatics 19, 2502–-2504 (2003).
29. Huang, X. et al. Overexpression of human reticulon 3 (hRTN3) in astrocytoma.
Clin. Neuropathol. 23, 1–-7 (2004).
30. Johnson, J. M. et al. Genome-wide survey of human alternative pre-mRNA
splicing with exon junction microarrays. Science 302, 2141–-2144 (2003).
Supplementary Information is linked to the online version of the paper at
www.nature.com/nature.
Acknowledgements This paper is dedicated to the memory of Stan Korsmeyer.
We thank members of the Vidal laboratory and the participants of the ORFeome
Meeting for discussions; the sequencing staff at Agencourt Biosciences for
technical assistance; E. Smith for his help with the figures; C. McCowan, A. Bird,
T. Clingingsmith and C. You for administrative assistance; and E. Benz,
S. Korsmeyer, D. Livingston, P. McCue, J. Song, B. Rollins and the DFCI Strategic
Planning Initiative for support. Our human interactome project is supported by
the DFCI High-Tech Fund (S. Korsmeyer), an Ellison Foundation grant awarded
to M.V., an NIH/NCI grant awarded to S. Korsmeyer, S. Orkin, G. Gilliland and
M.V., an ‘interactome mapping’ grant from NIH/NHGRI and NIH/NIGMS
awarded to F.P.R. and M.V., and a W.M. Keck Foundation grant awarded to
E. Benz, J. Marto, F.P.R. and M.V. Other support includes Taplin Funds for
Discovery (F.P.R., F.D.G. and G.F.B), a 2003 NSF Fellowship (D.S.G) and funding
from the Fonds National de la Recherche Scientifique, Belgium (M.D.).
Author Contributions Experiments and data analyses were coordinated by J.F.R.,
T.H. and K.V. High-throughput ORF cloning and yeast two-hybrid screens were
performed by J.F.R., T.H.K., A.D., N.L., N.A.G., J.R. and J.L. J.F.R developed the
high-throughput yeast two-hybrid strategy. Computational analyses were
performed by T.H., K.V., G.F.B., F.D.G., N.K., P.L., D.S.G., L.V.Z., S.L.W. and G.F.
Co-affinity purification experiments were performed by M.D., C.S., J.F.R., S.M.,
M.B., S.L. and J.S.A. C.F., E.L., S.C. and C.B. provided laboratory support. R.S.S.,
J.V., H.Y.Z., A.S. and M.E.C. helped with the overall interpretation of the data.
DNA sequencing was performed by S.B., R.S. and L.D.S. The manuscript was
written by J.F.R., K.V., M.E.C., D.E.H., F.P.R. and M.V. The project was conceived
by M.V. and co-directed by D.E.H., F.P.R. and M.V.
Author Information Reprints and permissions information is available at
npg.nature.com/reprintsandpermissions. The authors declare no competing
financial interests. Correspondence and requests for materials should be
addressed to M.V. (marc\_vidal@dfci.harvard.edu), F.P.R.
(fritz\_roth@hms.harvard.edu) or D.E.H. (david\_hill@dfci.harvard.edu).
LETTERS
NATURE|Vol 437|20 October 2005
1178
\end{document}
